\section{Background and related work}

\subsection{Information, selective prediction, and answerability}

Blackwell's comparison of statistical experiments formalizes when one
observation can reproduce another decision problem by discarding or garbling
information \citep{blackwell1951,blackwell1953}. The nested-channel
monotonicity in this report is a direct policy-set specialization of that
established idea. It does not imply that a fitted rule benefits from more
features, that a finite sample achieves the population frontier, or that
comparative records are trustworthy in a real monitoring network.

Selective classification and reject-option learning already study the
risk--coverage trade-off, optimal coverage subject to a desired risk, and the
cost of abstention \citep{chow1970,elyaniv2010,bartlett2008,geifman2017,franc2023}.
Hierarchical selective classification further permits a less-specific answer
when a fine label is not justified \citep{goren2024}. These are important
precedents, not contributions claimed here. MetaShift-Bench specializes the
object being selected: the answer is a synthetic local-versus-shared
\emph{scope} judgment under explicitly contrasting information channels.
The held-out envelope is deliberately descriptive. Calibration-selected
confidence does not itself provide a distribution-free conditional-risk
guarantee \citep{barberlimits2021}.

\subsection{Common, idiosyncratic, and weakly labelled changes}

Multivariate change-point work distinguishes common and idiosyncratic
components or tests common breaks in panels \citep{bai1998,barigozzi2018}.
Those methods motivate the scope distinction, but they do not make an
availability-normalized donor-participation continuum or this benchmark's
target-fixed paired design standard. Conversely, weak-supervision research
studies inference with partial or proxy labels \citep{ratner2016,cauchois2024}.
The v0.5 scope labels are generator-defined gold labels, whereas a real AQS
Method Code remains only a weak metadata anchor. Neither form of evidence
identifies a physical mechanism.

\subsection{Residual signatures, fault diagnosis, and sensor placement}

Analytical redundancy, residual signatures, fault distinguishability, and
sensor placement are established in model-based diagnosis
\citep{chowwillsky1984,krysander2008,taiebat2017}. They require physical
relations, fault models, or sensor-network assumptions that this project does
not possess. The structural certificate is therefore not claimed as a new
fault-diagnosis method, a sensor-observability result, or a real-device
isolation procedure. It is a bounded sufficient condition for abstention in
the stipulated synthetic score construction.

\subsection{Monitoring comparability and public metadata}

Monitoring comparability is an empirical property rather than a label inferred
from one administrative field. EPA publishes AQS file formats, Method Code and
POC references, and continuous-monitor comparability resources
\citep{epa_airdata_formats,epa_method_code,epa_poc,epa_pm25_comparability}.
Collocated Federal Reference Method and Federal Equivalent Method studies show
why configuration-specific comparison matters \citep{khan2024}. They do not,
however, provide dated physical-change ground truth for every historical AQS
Method Code transition.

The climate-record homogenization literature offers a useful structural
analogy: pairwise comparisons and station histories can expose
inhomogeneities, while the cause remains conditional on the available
documentation \citep{menne2009}. Air-pollution analyses similarly use
meteorological normalization to distinguish interventions or trends from
changing weather \citep{gagliardi2022,grange2019,zheng2023}. Low-cost PM2.5
sensor studies further motivate calibration and reference comparison
\citep{clements2017,barkjohn2021,chu2020}. Their instruments, target
populations, and calibration goals differ from a retrospective audit of
regulatory-monitor metadata. MetaShift-Bench uses this literature as
measurement-quality context, not as evidence that a particular AQS transition
has a particular physical cause.

\subsection{Reference series, counterfactuals, and change points}

Single-series change-point procedures include cumulative-sum monitoring,
segmentation methods such as PELT, and multiscale approaches
\citep{killick2012,truong2020,fryzlewicz2014}. These methods are valuable
baselines because they find patterns in one sequence; they cannot by
themselves determine whether a break is monitor-local or shared across a
region. The difference motivates matched regional variants in this
benchmark.

Synthetic control provides donor-weighted counterfactuals and in-space placebo
logic \citep{abadie2010,abadie2015}. Generalized, augmented, and synthetic
difference-in-differences extensions address other panel structures and
regularization choices \citep{xu2017,benmichael2021,arkhangelsky2021}.
The modern synthetic-control overview emphasizes that data requirements and
failure conditions remain substantive \citep{abadie2021}.
Modern difference-in-differences work also emphasizes assumptions and
heterogeneity under varied treatment timing \citep{callaway2021,goodmanbacon2021}.
MetaShift-Bench adopts the computational ideas of pre-event fitting and
reference comparison, but does not claim that the causal identification
assumptions of a policy study apply to a reported monitoring-method transition.

\subsection{Uncertainty, selection, and abstention}

Time-series resampling must respect serial dependence; circular moving-block
bootstrap methods are a natural descriptive tool for this purpose
\citep{kuensch1989}. Distribution-free predictive and conformal methods offer
another way to assess interval behavior when exchangeability assumptions are
made explicit \citep{vovk2005,lei2018,romano2019,barber2021}. In this report,
these methods motivate synthetic coverage evaluation rather than a claim that
real monitoring events enjoy generic calibration.

The evidence-tier task is also related to multiple testing, selection, and
abstention. Benjamini--Hochberg adjustment organizes an exploratory family of
screening quantities \citep{benjamini1995}; it does not turn those quantities
into causal probabilities. Post-selection inference highlights why selection
can invalidate naive uncertainty statements \citep{berk2013}. Research on
noisy labels underscores that observed proxies are not ground truth
\citep{liu2016,frenay2014}. MetaShift-Bench consequently records an
inconclusive tier when required comparative evidence is unavailable instead of
forcing a positive or negative physical diagnosis.

\subsection{Contribution boundary}

Table~\ref{tab:related-work-boundary} locates the benchmark relative to its
component literatures. The comparison concerns the scope of the reported
workflow, not a claim that MetaShift-Bench improves the underlying methods.

\begin{table}[tbp]
\centering
\caption{Relationship between MetaShift-Bench and contributing research
traditions.}
\label{tab:related-work-boundary}
\small
\begin{tabularx}{\linewidth}{@{}p{0.22\linewidth}X X@{}}
\toprule
Research tradition & Typical object and evidence & MetaShift-Bench boundary \\
\midrule
Synthetic control and DiD \citep{abadie2010,arkhangelsky2021} &
Pre-intervention panel outcomes and comparison units for an estimand under
explicit assumptions. &
Uses pre-anchor donor fitting and placebos as diagnostics; does not assert
policy-style causal identification for metadata transitions. \\
\addlinespace
Single-series change points \citep{killick2012,truong2020} &
One sequence and a change statistic or estimated break. &
Retains these methods as baselines but separates a target-local pattern from a
matched regional movement only when a qualified reference is available. \\
\addlinespace
Monitoring comparability and collocation \citep{khan2024,epa_pm25_comparability} &
Configuration-specific technical comparisons, collocation, and station
documentation. &
Treats auxiliary POC, QA, and document records as contextual evidence; no
metadata field is promoted to physical-device ground truth. \\
\addlinespace
Climate-record homogenization \citep{menne2009} &
Station histories plus pairwise evidence for record inhomogeneities. &
Adopts the audit logic of preserving candidate discontinuities and requiring
records for mechanism; it does not claim a homogenized corrected PM2.5 series. \\
\bottomrule
\end{tabularx}
\end{table}

The bounded joint contribution is a target-fixed donor-participation continuum
that evaluates channel-specific selective scope answerability and a
bounded-error structural abstention condition in a synthetic audit setting.
It combines established ingredients rather than claiming their invention:
information refinement, selective risk--coverage, common/idiosyncratic
changes, donor-weighted comparisons, pairwise monitoring comparisons, and
residual diagnosis are all prior work. Pairwise homogenization in particular
has long used reference comparisons and controlled simulations
\citep{menne2009,williams2012}. We did not identify, in a focused
primary-source audit, a source stating this complete joint construction. That
absence is not proof of priority. The report makes no ``first ever'' claim,
no generic risk--coverage claim, no estimator-superiority claim, and no real
monitoring-mechanism identification claim.
